\documentclass[a4paper,11pt]{ltjsarticle}

% --- パッケージの読み込み ---
\usepackage{amsmath, amssymb, amsthm, mathrsfs}
\usepackage{luatexja}
\usepackage{tikz}
\usetikzlibrary{cd, shapes, backgrounds, positioning, calc, arrows.meta, fit, patterns, decorations.pathreplacing}
\usepackage{hyperref}
\usepackage{xcolor}
\usepackage{listings}
\usepackage{tcolorbox}
\usepackage{booktabs}

% --- ハイパーリンク設定 ---
\hypersetup{
    colorlinks=true,
    linkcolor=blue!70!black,
    citecolor=blue!70!black,
    urlcolor=blue!70!black
}

% --- 定理環境の定義 ---
\theoremstyle{definition}
\newtheorem{definition}{定義}[section]
\newtheorem{theorem}[definition]{定理}
\newtheorem{lemma}[definition]{補題}
\newtheorem{example}[definition]{例}
\newtheorem{remark}[definition]{注釈}

% --- マクロ定義 ---
\newcommand{\N}{\mathbb{N}}
\newcommand{\R}{\mathbb{R}}
\newcommand{\OmegaSpace}{\Omega}
\newcommand{\F}{\mathcal{F}}
\newcommand{\Prob}{\mathbb{P}}
\newcommand{\Expect}{\mathbb{E}}
\newcommand{\CondExp}[2]{\Expect[#1 \mid #2]}
\newcommand{\Tau}{\tau}
\newcommand{\Tower}{\mathcal{T}}
\newcommand{\Layer}[1]{\mathrm{Layer}_{#1}}
\newcommand{\MinLayer}{\mathrm{minLayer}}

% --- タイトル情報 ---
\title{\textbf{構造塔による確率論の再構築：第1版フィルトレーション} \\ \large --- Bourbaki形式主義からマルチンゲールまで ---}
\author{
    Lean Code Generation: \textbf{CODEX (OpenAI)} \\
    TeX Explanation \& Typesetting: \textbf{Gemini (Google)}
}
\date{\today}

\begin{document}

\maketitle

\begin{abstract}
本稿では、一連のLean 4ファイルに基づき、数学的構造の階層性を抽象化した「構造塔 (Structure Tower)」の理論と、それを基盤とした確率論（特に停止時間とマルチンゲール）の形式化について解説する。Bourbaki流の一般構造から出発し、圏論的な普遍性を持つ構造塔を定義、それを$\sigma$-代数に適用することでフィルトレーションを再解釈し、最終的にマルチンゲールの理論へと接続する一連の流れを体系的に記述する。
\end{abstract}

\tableofcontents
\newpage

% =========================================================================
% Section 1: Bourbaki_Lean_Guide
% =========================================================================
\section{Bourbaki流のアプローチと構造の階層}
\label{sec:bourbaki}
\textit{Based on: \texttt{Bourbaki\_Lean\_Guide.md}}

本プロジェクトの出発点は、ニコラ・ブルバキの『数学原論』に見られる「母構造」の思想を、現代の定理証明支援系Leanで表現することにある。

\subsection{具体的な対象から抽象的な構造へ}
数学的対象を個別に扱うのではなく、それらが属する「順序集合 (Poset)」や「階層構造」として捉えることが推奨される。例えば、自然数の初期部分集合 $\{k \in \N \mid k \le n\}$ の族は、単なる集合の集まりではなく、包含関係によって順序付けられた構造をなす。

\subsection{構造塔 (Structure Tower) の萌芽}
このガイドでは、後の理論の核となる \texttt{StructureTower} の原型が提示されている。これは、順序集合 $I$ で添字付けられた集合族 $(\Layer{i})_{i \in I}$ であり、単調性（$i \le j \implies \Layer{i} \subseteq \Layer{j}$）を満たすものである。

\begin{figure}[h]
\centering
\begin{tikzpicture}[scale=1]
    \draw[fill=gray!10, draw=gray!50] (0,0) circle (2.5cm);
    \node at (0, 2.2) {全体空間 $X$};
    
    \draw[fill=blue!20, draw=blue!80] (0,0) circle (1.5cm);
    \node[blue!80] at (0, 1.2) {$\Layer{j}$};
    
    \draw[fill=red!20, draw=red!80] (0,0) circle (0.8cm);
    \node[red!80] at (0, 0) {$\Layer{i}$};
    
    \draw[->, thick] (0.5, 0.2) -- (1.2, 0.2);
    \node at (2.5, 0) {$i \le j$};
\end{tikzpicture}
\caption{構造塔の基本イメージ。順序に従って層が拡大していく。}
\end{figure}

% =========================================================================
% Section 2: CAT2_complete
% =========================================================================
\section{構造塔の圏論的定式化}
\label{sec:cat2}
\textit{Based on: \texttt{CAT2\_complete.md}}

前節のアイデアを厳密化し、圏論的な性質（普遍性、直積など）を持つ対象として \texttt{StructureTowerWithMin} を定義する。

\subsection{最小層を持つ構造塔}
単なる単調な集合族に加え、「各要素が属する最小の層」が一意に定まることを要請する。これにより、構造塔の射の一意性が保証され、圏論的な取り扱いが可能になる。

\begin{definition}[StructureTowerWithMin]
構造塔とは、以下の要素からなる組である。
\begin{itemize}
    \item 基礎集合 $X$ (carrier)
    \item 添字順序集合 $I$ (Index)
    \item 層構造 $\Layer{\cdot}: I \to \text{Set } X$
    \item \textbf{最小層関数} $\MinLayer: X \to I$
\end{itemize}
これらは以下の公理を満たす。
\begin{enumerate}
    \item \textbf{単調性}: $i \le j \implies \Layer{i} \subseteq \Layer{j}$
    \item \textbf{被覆性}: 任意の $x$ はある層に含まれる。
    \item \textbf{最小性}: $x \in \Layer{i} \iff \MinLayer(x) \le i$
\end{enumerate}
\end{definition}

\subsection{構造塔の射と圏}
構造塔の間の射 $(f, \phi): \Tower \to \Tower'$ は、基礎集合の写像 $f: X \to X'$ と添字の写像 $\phi: I \to I'$ の対であり、特に「最小層を保存する」ことが重要である。
\[ \MinLayer'(f(x)) = \phi(\MinLayer(x)) \]
この条件により、自由構造塔からの射の一意性（普遍性）が保証され、構造塔の全体は圏をなす。

% =========================================================================
% Section 3: SigmaAlgebraTower
% =========================================================================
\section{$\sigma$-代数とフィルトレーションの構造}
\label{sec:sigma_algebra}
\textit{Based on: \texttt{SigmaAlgebraTower.md}}

抽象的な構造塔理論を、測度論・確率論の基礎である $\sigma$-代数に応用する。ここからが確率論への具体的な展開の始まりである。

\subsection{$\sigma$-代数の構造塔}
標本空間 $\OmegaSpace$ の部分集合 $A$ に対して、「$A$ を可測にする最小の $\sigma$-代数」を考えることができる。これは $\sigma(A)$ と書かれ、構造塔における $\MinLayer(A)$ に相当する。
\[ \Layer{\F} = \{A \subseteq \OmegaSpace \mid A \in \F\} \]
ここで添字集合 $I$ は $\sigma$-代数全体の集合（包含関係による順序）である。

\subsection{フィルトレーション (Filtration)}
確率論におけるフィルトレーション $(\F_n)_{n \in \N}$ は、情報の増大を表す $\sigma$-代数の列である。
\[ \F_0 \subseteq \F_1 \subseteq \F_2 \subseteq \dots \]
これは、添字集合を自然数 $\N$ とした構造塔の特別な場合（または構造塔への入力）と見なせる。
`FiltrationToTower` 関数により、フィルトレーションから「可測事象の階層」としての構造塔が構成される。

\begin{lemma}[情報の単調性]
フィルトレーションにおいて、時刻 $n$ で可測な事象は、任意の未来 $m \ge n$ においても可測である。
\end{lemma}

% =========================================================================
% Section 4: StoppingTime_MinLayer
% =========================================================================
\section{停止時間と最小層の対応}
\label{sec:stopping_time}
\textit{Based on: \texttt{StoppingTime\_MinLayer.md}}

ここで、確率論の核心的な概念である「停止時間」が、構造塔理論の $\MinLayer$ として自然に解釈できることを示す。

\subsection{停止時間の定義}
停止時間 $\Tau: \OmegaSpace \to \N$ とは、任意の $n$ について $\{\omega \mid \Tau(\omega) \le n\} \in \F_n$ を満たす確率変数である。

\subsection{構造的解釈}
フィルトレーション $\F$ から誘導される構造塔 $\Tower_\F$ を考える。標本点 $\omega$ に対して、単集合 $\{\omega\}$ が初めて可測になる（識別可能になる）時刻を考えることができる（離散空間の場合）。
より一般的に、停止時間 $\Tau$ は、標本空間上の「各点が決定される時刻」を表すランク関数と見なせる。

\begin{theorem}[停止集合と層の一致]
停止時間 $\Tau$ が定める停止集合 $\{\Tau \le n\}$ は、構造塔の第 $n$ 層そのものである。
\[ \{\omega \mid \Tau(\omega) \le n\} = \Layer{n}(\Tower_\Tau) \]
また、$\MinLayer$ 関数は停止時間 $\Tau$ そのものに対応する。
\end{theorem}

\begin{figure}[h]
\centering
\begin{tikzpicture}[xscale=1.5, yscale=1]
    \draw[->] (0,0) -- (4,0) node[right] {時刻 $n$};
    
    % Filtration boxes
    \node[draw, fill=blue!10, minimum width=1cm, minimum height=1cm] at (0,1) {$\F_0$};
    \node[draw, fill=blue!20, minimum width=1.5cm, minimum height=1.2cm] at (1.2,1.1) {}; \node at (1.2,1.1) {$\F_1$};
    \node[draw, fill=blue!30, minimum width=2cm, minimum height=1.4cm] at (2.5,1.2) {}; \node at (2.5,1.2) {$\F_2$};
    
    % Stopping Time
    \node[red] at (2, -0.5) {$\Tau(\omega) = 2$};
    \draw[red, ->, dashed] (2, -0.2) -- (2, 0.8);
    \node[align=left, font=\footnotesize] at (2, 2.5) {$\omega$ が識別可能に\\なる最小の層};
\end{tikzpicture}
\caption{フィルトレーションと停止時間。時間は構造塔の「深さ」に対応する。}
\end{figure}

\subsection{停止過程 (Stopped Process)}
確率過程 $X$ を停止時間 $\Tau$ で止めた過程 $X^\Tau$ を定義する。
\[ X^\Tau_n(\omega) = X_{\min(n, \Tau(\omega))}(\omega) \]
これは構造塔上での「値の凍結」操作として定式化される。

% =========================================================================
% Section 5: Martingale_StructureTower
% =========================================================================
\section{マルチンゲールと構造的保存則}
\label{sec:martingale}
\textit{Based on: \texttt{Martingale\_StructureTower.md}}

最後に、この構造塔の上で定義されるマルチンゲールについて解説する。

\subsection{条件付き期待値とマルチンゲール}
条件付き期待値 $\CondExp{\cdot}{\F_n}$ は、構造塔の第 $n$ 層への射影（平滑化）と見なせる。
マルチンゲール $M = (M_n)$ は、この射影に関して不変な過程である。
\[ \CondExp{M_{n+1}}{\F_n} = M_n \]

\subsection{構造塔上のマルチンゲール定義}
Leanファイルでは、以下の要素を持つ構造体としてマルチンゲールを定義している。
\begin{itemize}
    \item \texttt{filtration}: 基礎となる構造塔（フィルトレーション）
    \item \texttt{process}: 適合する確率過程
    \item \texttt{adapted}: 各時刻での強可測性
    \item \texttt{integrable}: 可積分性
    \item \texttt{martingale}: マルチンゲール等式
\end{itemize}

\subsection{停止マルチンゲール}
停止時間 $\Tau$ で止めたマルチンゲール $M^\Tau$ もまたマルチンゲールになる（有界な場合など）。これは、構造塔の射（停止操作）がマルチンゲールという「構造」を保存することを意味している。この事実は \texttt{stoppedProcess\_martingale\_property\_of\_bdd} として形式化されている。

\section{結論}
本稿で解説した一連の形式化は、確率論を「解析的な計算」の対象から、「順序構造上の代数的な操作」の対象へと昇華させる試みである。
\begin{enumerate}
    \item \textbf{Bourbaki Guide}: 構造の階層化という哲学。
    \item \textbf{CAT2}: 最小層を持つ構造塔の厳密な定義。
    \item \textbf{SigmaAlgebra}: 確率論の舞台の構造化。
    \item \textbf{StoppingTime}: 確率的時間の構造的解釈。
    \item \textbf{Martingale}: 構造を保存するダイナミクス。
\end{enumerate}
このアプローチにより、任意停止定理などの深い結果が、構造の普遍性として自然に理解される道が開かれる。

\end{document}